%%=============================================================================
%% Inleiding
%%=============================================================================

\chapter{\IfLanguageName{dutch}{Inleiding}{Introduction}}%
\label{ch:inleiding}

In onderzoeksprojecten worden berekeningen vaak eerst uitgewerkt door onderzoekers en pas daarna omgezet naar software. Dat is logisch: onderzoekers kennen het domein en de formules, terwijl ontwikkelaars instaan voor de technische uitwerking. Toch kan die verdeling ook voor problemen zorgen. Deze bachelorproef onderzoekt dat probleem binnen de context van ILVO en het EMAV-project.                                                                                   
Bij ILVO (Instituut voor Landbouw-, Visserij- en Voedingsonderzoek) worden rekenmodellen gebruikt om onder andere milieu-impact en emissies binnen de landbouwsector te berekenen. Zulke modellen vertrekken vanuit wetenschappelijke kennis, formules en aannames die door onderzoekers worden opgesteld. Wanneer een model in een toepassing gebruikt moet worden, volstaat een Excel-bestand of los document meestal niet meer.                                                 
De berekeningen moeten dan worden omgezet naar software die betrouwbaar, onderhoudbaar en bruikbaar is voor eindgebruikers. In de onderzochte context gebeurt dat binnen een .NET-omgeving, met C\# als programmeertaal.                                 De onderzoekers die met de modellen werken, zijn niet noodzakelijk vertrouwd met C\#. Ze zijn vaak wel vertrouwd met de achterliggende formules, de parameters en de betekenis van de resultaten. Daardoor ontstaat een praktische vraag: hoe kan de software zo worden ingericht dat onderzoekers de berekeningen beter kunnen volgen en eventueel zelf kleine experimenten kunnen uitvoeren?                                                                                     
In de praktijk ontstaat er een afstand tussen de wetenschappelijke versie van het model en de technische implementatie ervan. Onderzoekers leveren formules, parameters of voorbeeldberekeningen aan. Ontwikkelaars vertalen die vervolgens naar C\#-code, koppelen ze aan databanken en verwerken ze in de rest van de applicatie.                                                                                                                                                
Dat werkt zolang het model stabiel blijft. Zodra een onderzoeker een formule wil aanpassen, een parameter wil testen of een alternatief scenario wil vergelijken, is er opnieuw technische hulp nodig. De onderzoeker kan de berekening meestal niet rechtstreeks uitvoeren in de softwarecontext waarin het model echt gebruikt wordt. Omgekeerd is het voor ontwikkelaars niet altijd eenvoudig om de wetenschappelijke bedoeling achter elke berekeningsstap te blijven volgen. 
De kern van het probleem is dat de berekeningslogica na verloop van tijd moeilijk zichtbaar wordt voor de onderzoeker. De code voert het model wel uit, maar de wetenschappelijke redenering zit verspreid over documentatie, broncode, overleg en voorbeeldbestanden. Daardoor kan de toepassing voor onderzoekers aanvoelen als een black box.                                                                                                                                   
Een bijkomend probleem is dat documentatie snel achterloopt op de implementatie. Wanneer code wijzigt maar de uitleg niet mee aangepast wordt, ontstaat "documentation drift". Dat maakt het moeilijker om na te gaan of de software nog overeenkomt met het model zoals onderzoekers het bedoelen. Deze bachelorproef onderzoekt daarom hoe code, documentatie en experimenteerruimte dichter bij elkaar gebracht kunnen worden.                                                  
Deze bachelorproef probeert niet het volledige EMAV-project te herwerken. Ook het volledige emissiemodel zelf wordt niet opnieuw ontworpen. De scope ligt bij een proof-of-concept waarin wordt nagegaan hoe C\#-logica beschikbaar gemaakt kan worden in een Python-omgeving en hoe een alternatieve implementatie getest kan worden zonder de hoofdapplicatie rechtstreeks aan te passen.                                                                                        
Vanuit deze problematiek is de volgende centrale onderzoeksvraag geformuleerd:
\textit{Hoe kunnen we de kloof tussen IT-logica en wetenschappelijke analyse verkleinen, zodat onderzoekers zonder diepgaande programmeerkennis inzicht krijgen in hun modellen en er zelf mee kunnen experimenteren?}                                   Het doel van dit onderzoek is om een haalbare werkwijze te verkennen waarmee onderzoekers meer inzicht krijgen in de berekeningen, zonder dat de bestaande software volledig herbouwd moet worden. Daarvoor wordt een proof-of-concept uitgewerkt. De focus ligt niet alleen op documentatie, maar ook op de vraag hoe onderzoekers op een gecontroleerde manier met de berekeningen kunnen experimenteren.

Eerst wordt in de stand van zaken bekeken welke bestaande technieken en werkwijzen relevant zijn, zoals Python.NET, computationele notebooks, literate programming en containerisatie. 
Daarna licht de methodologie toe hoe het probleem en de mogelijke oplossing onderzocht werden.
Dan volgen hoofdstukken over de alternatieve vormen van documentatie, de concepten en technieken van clean code en documentation drift.
Hierop volgt een hoofdstuk over de implementatie van python in C\# en hoe de complexiteit sterk veranderd gebaseerd op hoeveel in python gewenst is.

